\documentclass[11pt,a4paper]{article}
\usepackage[utf8]{inputenc}
\usepackage{amsmath,amssymb}
\usepackage{graphicx}
\usepackage{booktabs}
\usepackage{hyperref}
\usepackage{geometry}
\usepackage{natbib}
\geometry{margin=1in}

\title{Cultural Alignment of Language Models through\\Figurative Language: Continue Pretraining}

\author{}
\date{}

\begin{document}
\maketitle

\begin{abstract}
We present a pipeline for improving the cultural understanding of language models through figurative language. Our approach currently targets English and Chinese, with planned extensions to Hindi and Arabic, by leveraging large-scale idiom datasets ($\sim$21K English, $\sim$28K Chinese idioms with figurative and literal meanings). We filter and deduplicate the C4/mC4 corpus to extract high-quality documents containing culturally meaningful idioms, cap per-idiom document frequency to ensure broad coverage, and augment documents with idiom metadata tags exposing literal and figurative meanings. We describe the data construction methodology, discuss our plans for extending the pipeline to Hindi and Arabic, and outline our evaluation plan on cultural benchmarks.
\end{abstract}

\section{Introduction}

Language models trained on large web corpora acquire broad linguistic competence but often lack deep understanding of culture-specific behaviors, norms, and values. Figurative language---idioms, proverbs, and metaphorical expressions---encodes cultural knowledge in ways that literal text does not. An idiom's figurative meaning frequently reflects historical attitudes, social norms, and philosophical worldviews unique to its culture of origin. For instance, the Chinese chengyu ``pao zhuan yin yu'' (throw a brick to attract jade) encodes a Confucian ethic of humility and collective improvement, while the English idiom ``pull yourself up by your bootstraps'' reflects an Anglo-American emphasis on individual self-reliance. These figurative meanings are cultural artifacts, and a model that understands them gains access to a layer of cultural reasoning that surface-level language competence alone cannot provide.

We propose a continue pretraining pipeline that systematically injects such cultural knowledge into language models. The pipeline exposes the model to a large volume of natural text containing idioms in context, augmented with explicit meaning annotations. Our pipeline operates independently per language---English prompts and idioms for English culture, Chinese prompts and chengyu for Chinese culture---to avoid cross-cultural contamination and ensure each language's cultural knowledge is learned on its own terms. While our current implementation covers English and Chinese, the pipeline is designed to be language-agnostic, and we plan to extend it to Hindi and Arabic, two languages with rich figurative traditions and distinct cultural value systems.

\section{Data Resources}

Our pipeline draws on two primary data sources. The first is a pair of comprehensive idiom datasets: approximately 21K English idioms with figurative meanings, sourced from multiple dictionaries and filtered to retain only entries with non-trivial figurative content, and approximately 28K Chinese chengyu with figurative meanings and classical source citations. Each idiom entry contains the idiom text, a list of key conceptual entities, literal meanings (often with classical source quotations for Chinese), and figurative meanings. Trivial English entries---single common words, basic grammatical constructions, and phrases understood literally without cultural context---are removed using a combination of heuristic filtering (requiring non-empty figurative meanings) and LLM-based classification.

The second data source is the C4 corpus~\citep{raffel2020exploring} (English: $\sim$365M documents) and its multilingual counterpart mC4 (Chinese subset), which serve as the raw material for continue pretraining.

\section{Continue Pretraining}

The goal of continue pretraining is to expose the model to a large volume of natural text in which idioms appear in authentic contexts, while ensuring the model can associate each idiom with its figurative meaning. This stage involves three components: document filtering, deduplication with frequency capping, and idiom metadata tagging.

\paragraph{Document filtering.} We filter the C4/mC4 corpus to extract documents containing at least one idiom from our dataset. We stream the full corpus and use Aho-Corasick multi-pattern matching to test each document against all $\sim$21K English or $\sim$28K Chinese idiom patterns simultaneously in a single pass, with case-insensitive matching for English. Given the scale of the English C4 corpus ($\sim$365M documents), we partition the workload across multiple SLURM jobs using document offset resumption, achieving a processing rate of approximately 50K documents per minute. During filtering, we also accumulate per-idiom document counts that record how many documents contain each idiom, which feed into the subsequent deduplication and frequency capping stage. Additionally, our pipeline supports an infini-gram~\citep{liu2024infinigram} mode that indexes the filtered corpus and enables efficient per-idiom frequency counting and targeted document extraction without re-scanning the full dataset.

\paragraph{Deduplication and frequency capping.} A naive filtering approach produces a heavily skewed distribution: common idioms containing high-frequency words (e.g., ``hand,'' ``heart,'' ``time'' in English, or ``xin'' (heart), ``ren'' (person), ``tian'' (sky) in Chinese) appear in hundreds of thousands of documents, while rare idioms may appear in only a handful. Training on this skewed distribution would cause the model to over-attend to common idioms while neglecting rare, culturally specific ones. To address this, we perform per-idiom frequency capping at a maximum of 10,000 documents per idiom. When an idiom appears in more than 10,000 documents, we select a high-quality subset by prioritizing documents where the idiom appears in meaningful context (not in lists, dictionary entries, or metadata), the document is of sufficient length to provide contextual richness, and the document is not a near-duplicate of other selected documents. Near-duplicate detection uses hash-based comparison on the first 500 characters. By capping frequent idioms while retaining all documents for rare ones, we ensure broad and balanced coverage across the full idiom vocabulary.

\paragraph{Idiom metadata tagging.} A key component of our continue pretraining approach is augmenting filtered documents with structured idiom metadata. When a document contains one or more idioms from our dataset, we append a knowledge block to the document text that lists each detected idiom along with its figurative and literal meanings from our idiom dataset. For Chinese chengyu, this includes classical source citations. This tagging creates explicit meaning associations: the model learns to connect each idiom with its figurative meaning through co-occurrence in the training data, rather than relying solely on contextual inference.

\section{Instruction Tuning}

We fine-tune the model on instruction-following data that requires idiom understanding and cultural reasoning.
We generate the instruction-following data through the following ways:
1. Based on the comprehensive idiom datasets for each language, we generate a set of training data that generate questions about each idiom and ask the model to answer with some idioms that are related to the question and provide its reasoning. We will ask an LLM to generate contexts that are related to each idiom's figurative meaning. For example, Given an english idiom "throw a brick to attract jade", we will ask the model to generate a context that is related to the figurative meaning of the idiom: xxxx, covering what would be the appropriate context/behavior/situation that we could describe it using the idiom. In the output, we will ask the model to answer the question using this idiom, and provide the reasoning for why this idiom is the most appropriate one to use in the context.
2. We will link several idioms that share similar meanings together, and provide generative questions, asking the model to find idioms that share teh similar meaning to one of the idioms, and ask the model to provide the reasonign for why their meanings are similar.
3. We will compare idioms in different languages. For example, we will find idioms in English and Chinese that are related to the same figurative meaning. We will provide one idiom in English, and ask the model to provide one or more Chinese idioms that are related to the same concept, and (1) provide the reasoning for why their meanings are similar. We will also provide the English and Chinese idiom together, asking the model to analyze the cultural differences between the two languages if there are any.
4. For highly frequent entities in each language's idiom dataset, we will generate questions that require the model to understand the cultural significance of the entity. For example, for the entity "heart" in English, we will generate questions that require the model to understand the cultural significance of the entity, such as "What is the cultural significance of the heart in English culture?", "What is the cultural significance of the heart in Chinese culture?", etc. (2) We will also compare the cultural significance of the entity in different languages. (3 and 4 differs in the sense of figurative menaing vs. used entity in the idiom)

\section{Evaluation Plan}

We plan to evaluate the trained model on several cultural benchmarks to assess whether the two-stage pipeline improves culture-related understanding and reasoning. These include cultural natural language inference tasks requiring cultural knowledge for correct entailment judgments; the World Value Survey, which tests alignment with cultural value distributions across different populations; and direct evaluation of figurative meaning comprehension in both languages, measuring whether the model can correctly identify and explain the cultural significance of idioms encountered in novel contexts.

\section{Extension to Hindi and Arabic}

While our current implementation focuses on English and Chinese, the pipeline is designed to generalize to additional languages. We plan to extend coverage to Hindi and Arabic, two languages with rich figurative traditions and cultural contexts that are underrepresented in current language model training. Hindi possesses a wealth of muhavare (idioms) and lokoktiyan (proverbs) that encode South Asian cultural values around family hierarchy, dharma, and social obligation, while Arabic amthal (proverbs) and kinayat (figurative expressions) reflect cultural norms rooted in Islamic tradition, tribal honor codes, and the centrality of hospitality.

For both languages, the extension follows the same methodology. First, we will compile idiom datasets with figurative and literal meanings from established dictionaries and linguistic resources. Second, we will filter the corresponding mC4 language subsets using the same Aho-Corasick matching and per-idiom frequency capping strategy. The language-agnostic design of our pipeline means that adding a new language requires only the idiom dataset and the mC4 subset, with no changes to the core infrastructure.

\section{Conclusion}

We have presented a continue pretraining pipeline for cultural alignment of language models through figurative language. The pipeline provides broad, balanced exposure to idioms in natural context with explicit meaning annotations. By operating independently per language, capping per-idiom document frequency for balanced coverage, and tagging documents with structured idiom metadata, our approach provides a principled methodology for injecting cultural knowledge into language models. With English and Chinese as our initial targets and planned extensions to Hindi and Arabic, the pipeline aims to improve cultural understanding across a diverse set of linguistic and cultural traditions.

\bibliographystyle{plainnat}
\begin{thebibliography}{9}
\bibitem[Liu et al.(2024)]{liu2024infinigram}
Liu, J., Zhu, Z., et al.
\newblock Infini-gram: Scaling Unbounded n-gram Language Models to a Trillion Tokens.
\newblock \textit{arXiv preprint arXiv:2401.02067}, 2024.

\bibitem[Raffel et al.(2020)]{raffel2020exploring}
Raffel, C., Shazeer, N., Roberts, A., Lee, K., Narang, S., Matena, M., Zhou, Y., Li, W., \& Liu, P.~J.
\newblock Exploring the limits of transfer learning with a unified text-to-text transformer.
\newblock \textit{Journal of Machine Learning Research}, 21(140):1--67, 2020.
\end{thebibliography}

\end{document}
